% !TeX spellcheck = hu_HU
% !TeX encoding = UTF-8
% !TeX program = xelatex
%----------------------------------------------------------------------------
\section{Különböző nyelvek együttes használata}\label{sec:uses:mix}
%----------------------------------------------------------------------------

A C4 nyelvnek, amellett, hogy egy-egy \gls{DSL} megvalósítását egyszerű megadni benne, van egy másik kiemelkedő erőssége.
Mivel a benne definiált nyelvek \gls{eDSL}-ek a C4 nyelven belül, így ezeket a nyelveket lehetséges olyan módon kombinálni, ahogy az a fejlesztő számára a legkényelmesebb megoldást tudja nyúltani az adott feladathoz.
Ez igaz bármilyen két beágyazott nyelv esetén, amennyiben azok nem definiálnak olyan függvényeket, amelyek névütközést okoznak -- erre a problémára is nyújtana megoldást \aref{sec:ftr:argfun}.~részben bemutatott továbbfejlesztési lehetőség.

Az alábbi példában egy parancssoros programot valósítok meg, amelyik a felhasználó által adott parancsok hatására egy állapotgépen belül vált állapotokat, majd a megfelelő lépések esetén a feladat ütemező által leírt részfeladatokat futtatja.
A használt \gls{DSL}-ek kódja azonos az előzőekben definiáltakkal, mind kettő megtalálható a függelékben.

Az alkalmazás kódja megtalálható \aref{fig:ex-mix}.~ábrán.
A cél egy demóként szolgálni egy olyan parancssoros rendszer felhasználói felületére, amelyikben az üzleti logikát -- adatok szűrése egy fájlban megadott kulcs-érték párok alapján, majd valamiféle riport előállítása -- egy feladatok hálójával lehetséges megadni.
A tényleges adatok szűrése nincs megvalósítva, mert a demó megértését csak bonyolítaná: a jelenleg cél csupán az, hogy bemutassam, milyen egyszerűen lehetséges két \gls{DSL}-t együttesen használni egy C4 programon belül, nem egy teljes alkalmazás fejlesztése.

Látszik a kódon, hogy a felhasználói felületnek az állapottartó jellegét jól meg lehet valósítani az állapotgépet felhasználva.
Ezzel párhuzamosan a futtatni szánt ``üzleti logikához'' tartozó kódot részekre lehet bontani, amelyet az ütemezőn keresztül lehet kényelmesen megvalósítani a függőségeket felhasználva.
A program ilyen módon való két részre bontásával a függőségek szétválasztása is automatikusan megtörténik.

\begin{figure}[htb]
    \begin{subfigure}{.46\linewidth}
        \begin{lstlisting}[numbers=left,gobble=12,morekeywords={runtime}]
            let runtime make
              rule "report" <: "run-filter": {
                println "report from filtered.csv"
              }
              rule "run-filter": {
                println "filter data.csv into filtered.csv"
              }

              rule "new-filter" <: "reset-filter"
                                 , "add-filter"
                                 : { }
              rule "reset-filter": {
                fprintln "filter.csv" "key;value"
              }
              rule "add-filter": {
                let key { println "Key: " readln }
                let value { println "Value: " readln }

                fappendln "filter.csv" strcat key
                                       strcat ";"
                                              value
              }
            done_make
        \end{lstlisting}
        \caption{Az ,,üzleti logika'' feladatok hálózatára bontott utasításokkal}
    \end{subfigure}
    \hfill
    \begin{subfigure}{.46\linewidth}
        \begin{lstlisting}[numbers=left,gobble=12,morekeywords={runtime}]
            let ui fsm
              start state "main-menu"
                enter { println "Welcome!" }
              final state "exit"
                enter { println "Bye!" }

              state "filter"
                enter { println "Entering filter edit mode" }
                leave { println "Leaving filter edit mode" }

              on "exit" from "main-menu" to "exit"

              on "generate" from "main-menu" to "main-menu"
                do { run runtime "report" }

              on "filter" from "main-menu" to "filter"
              on "exit" from "filter" to "main-menu"

              on "reset" from "filter" to "filter"
                do { run runtime "reset-filter" }
              on "set" from "filter" to "filter"
                do { run runtime "new-filter" }
              on "add" from "filter" to "filter"
                do { run runtime "add-filter" }
            done
        \end{lstlisting}
        \caption{A felhasználói felület megvalósítása állapotgéppel}
    \end{subfigure}
    \caption{Példa a C4 nyelvben megvalósított \gls{eDSL}-ek együttes használatára. A két részletet összekapcsoló „runtime” objektum a jobb átláthatóság kedvéért, ki van emelve.}\label{fig:ex-mix}
\end{figure}

Amit \aref{fig:ex-mix}.~ábrán érdemes megfigyelni, hogy a felhasználói felületet megvalósító állapotgép éleire írt akciók milyen könnyedén képesek egy teljesen független \gls{DSL}-ben megadott részét felhasználni a programnak.
Ahol az akciókban a {\tt run} függvény szerepel, ott a felhasználói felület belehív az üzleti logikába.

















